\section{L'Hospital's rule}
L'Hospital's rule compares a quotient with a quotient of derivatives; it does not say that every limit may be evaluated by differentiating. In the $0/0$ case, both numerator and denominator must tend to zero. Ross also permits the case in which the absolute value of the denominator tends to infinity, without requiring a numerator limit. In either case the functions must be differentiable on the relevant deleted neighborhood, the derivative of the denominator must be nonzero there, and the derivative quotient must have a limit. That limit may be finite or infinite.

For positive quantities with variable exponents, take a logarithm first. If $A(x)>0$, then $A(x)^{B(x)}=\exp(B(x)\log A(x))$; after finding the exponent's limit, use continuity of the exponential.

\exercise{30.1}{Four introductory limits}
\goal Evaluate each limit and justify the method, using the actual numerator and denominator in the printed exercise.
\step{(a) The quotient $(e^{2x}-\cos x)/x$ at zero.}
Both numerator and denominator tend to zero, they are differentiable, and the denominator derivative is $1\ne0$. L'Hospital's rule therefore gives
\[
 \lim_{x\to0}\frac{e^{2x}-\cos x}{x}
 =\lim_{x\to0}(2e^{2x}+\sin x)=2.
\]
\step{(b) The quotient $(1-\cos x)/x^2$ at zero.}
The original form is $0/0$, and the derivative of $x^2$ is nonzero for $x\ne0$. The derivative quotient has the known limit
\[
 \lim_{x\to0}\frac{\sin x}{2x}=\frac12.
\]
L'Hospital's rule thus gives $1/2$ for the original limit. The last step uses the basic sine limit, not another unjustified cancellation at zero.
\step{(c) The quotient $x^3/e^{2x}$ at positive infinity.}
At every differentiation stage the exponential denominator tends to infinity and has nonzero derivative. Successive derivative quotients reduce the polynomial degree:
\[
 \frac{x^3}{e^{2x}},\quad \frac{3x^2}{2e^{2x}},\quad
 \frac{6x}{4e^{2x}},\quad \frac6{8e^{2x}}.
\]
The final quotient tends to zero. Applying the rule backward through these three stages proves that the original limit is $0$.
\step{(d) A difference of square roots at zero.}
For $0<|x|<1$, multiply by the conjugate:
\[
 \frac{\sqrt{1+x}-\sqrt{1-x}}x
 =\frac{(1+x)-(1-x)}{x(\sqrt{1+x}+\sqrt{1-x})}
 =\frac2{\sqrt{1+x}+\sqrt{1-x}}\longrightarrow1.
\]
This direct algebraic method is shorter than L'Hospital's rule and also keeps the real domain explicit.
\answer{The four limits are $2$, $1/2$, $0$, and $1$.}

\exercise{30.2}{Cancellation, repeated differentiation, and a variable power}
\goal Evaluate the four limits, checking why each reduction is valid.
\step{(a) The quotient $x^3/(\sin x-x)$.}
The numerator and denominator and their first two derivative pairs tend to zero. Differentiating each pair produces
\[
 \frac{x^3}{\sin x-x},\quad
 \frac{3x^2}{\cos x-1},\quad
 \frac{6x}{-\sin x},\quad
 \frac6{-\cos x}.
\]
On a sufficiently small deleted neighborhood of zero, each denominator derivative used in passing to the next pair is nonzero: these are $\cos x-1$, $-\sin x$, and $-\cos x$. The final quotient tends to $-6$, so three valid applications of L'Hospital's rule yield $-6$.
\step{(b) The quotient $(\tan x-x)/x^3$.}
This is $0/0$ near zero, where tangent is defined. The denominator derivative $3x^2$ is nonzero for $x\ne0$. The derivative quotient is
\[
 \frac{\sec^2x-1}{3x^2}
 =\frac13\left(\frac{\tan x}x\right)^2\longrightarrow\frac13,
\]
because $\tan x/x=(\sin x/x)/\cos x\to1$. Thus the requested limit is $1/3$.
\step{(c) The difference $1/\sin x-1/x$.}
Combine the fractions before applying any limit rule:
\[
 \frac1{\sin x}-\frac1x=\frac{x-\sin x}{x\sin x}.
\]
The first derivative quotient is $(1-\cos x)/(\sin x+x\cos x)$, again of form $0/0$. Its derivative quotient is
\[
 \frac{\sin x}{2\cos x-x\sin x}\longrightarrow0.
\]
The final denominator tends to $2$ and hence is nonzero near zero. The intermediate denominator $\sin x+x\cos x$ is nonzero for sufficiently small nonzero $x$, because its ratio to $x$ tends to $2$. Thus both uses of L'Hospital's rule are legitimate and the original limit is $0$.
\step{(d) The expression $(\cos x)^{1/x^2}$.}
For $|x|<\pi/2$, the base is positive, so its logarithm is $\log(\cos x)/x^2$. This is $0/0$, and its derivative quotient is
\[
 \frac{-\tan x}{2x}\longrightarrow-\frac12.
\]
L'Hospital's rule gives the logarithmic limit $-1/2$. Continuity of the exponential now gives the original limit $e^{-1/2}$.
\answer{The answers are $-6$, $1/3$, $0$, and $e^{-1/2}$.}

\exercise{30.3}{Recognizing when differentiation is unnecessary}
\goal Evaluate the four limits, including a correction to a printed selected answer.
\step{(a) The quotient $(x-\sin x)/x$ at positive infinity.}
For the expression printed in the exercise,
\[
 \frac{x-\sin x}x=1-\frac{\sin x}x,
 \qquad\left|\frac{\sin x}x\right|\le\frac1x\longrightarrow0.
\]
The limit is therefore $1$. The selected-answer entry $0$ in the linked edition is inconsistent with this printed expression. Differentiating would give $1-\cos x$, which has no limit; this simply means L'Hospital's rule is not useful here. Its failure to supply a derivative-quotient limit says nothing against the elementary estimate above.
\step{(b) The power $x^{\sin(1/x)}$ at positive infinity.}
Its logarithm is
\[
 \sin(1/x)\log x
 =\frac{\sin(1/x)}{1/x}\,\frac{\log x}{x}.
\]
The first factor tends to $1$. The second tends to $0$ by L'Hospital's rule, since $(\log x)'/x'=1/x\to0$ and $x\to\infty$. Thus the logarithm tends to zero and the power tends to $e^0=1$.
\step{(c) The quotient $(1+\cos x)/(e^x-1)$ as $x\to0^+$.}
The numerator tends to $2$, while the denominator tends to zero through positive values. Consequently the limit is $+\infty$. Explicitly, near zero the numerator is greater than $1$. For any prescribed $M>0$, continuity of $e^x$ makes $0<e^x-1<1/M$ for sufficiently small $x>0$, giving the quotient greater than $M$. This is not a $0/0$ situation.
\step{(d) The quotient $(1-\cos 2x-2x^2)/x^4$ at zero.}
For transparency, list all derivatives through the first nonzero limiting values:
\[
\begin{array}{c|c|c}
 \text{order}&\text{numerator derivative}&\text{denominator derivative}\\ \hline
 0&1-\cos2x-2x^2&x^4\\
 1&2\sin2x-4x&4x^3\\
 2&4\cos2x-4&12x^2\\
 3&-8\sin2x&24x\\
 4&-16\cos2x&24
\end{array}
\]
At orders $0$ through $3$ both entries tend to zero. Every denominator derivative required by the rule is nonzero for nonzero $x$ near zero. The last ratio tends to $-16/24=-2/3$. Four applications, read backward, prove that this is the requested limit.
\answer{For the printed questions, the limits are $1$, $1$, $+\infty$, and $-2/3$.}

\exercise{30.4}{Changing a right-hand zero limit into an infinite limit}
\goal For $g(y)=f(1/y)$, prove equivalence of the two limits, including infinite limits.
\idea The substitution is reversible: $y\to+\infty$ corresponds exactly to $x=1/y\to0^+$.
\step{From a right-hand limit of $f$ to a limit of $g$.}
Suppose first $f(x)\to L\in\R$ as $x\to0^+$. Given $\varepsilon>0$, choose $\delta>0$, with $\delta<a$ if $a$ is finite, such that $0<x<\delta$ implies $|f(x)-L|<\varepsilon$. If $y>1/\delta$, then $0<1/y<\delta$, so $|g(y)-L|<\varepsilon$. This proves $g(y)\to L$ as $y\to\infty$.
\step{Reverse the argument.}
Suppose $g(y)\to L\in\R$. Choose $Y>0$ inside its domain such that $y>Y$ implies $|g(y)-L|<\varepsilon$. Let $\delta>0$ be at most $1/Y$ and small enough that $(0,\delta)\subset(0,a)$. Then $0<x<\delta$ implies $1/x>Y$, and
\[
 |f(x)-L|=|g(1/x)-L|<\varepsilon.
\]
Thus $f(x)\to L$ as $x\to0^+$.
\step{Use the same substitution for infinite limits.}
For a limit $+\infty$, replace $|f(x)-L|<\varepsilon$ and $|g(y)-L|<\varepsilon$ by $f(x)>M$ and $g(y)>M$, respectively, for arbitrary real thresholds $M$. The same reciprocal choices of $Y$ and $\delta$ give both implications. For $-\infty$, use the inequalities $f(x)<M$ and $g(y)<M$. No other part of the proof changes.
\answer{Either limit exists exactly when the other does, and their finite or infinite values are equal. The convention $a^{-1}=0$ when $a=\infty$ gives precisely the correct transformed domain.}

\exercise{30.5}{Three exponential limits}
\goal Evaluate the limits by taking logarithms on domains where the bases are positive.
\step{(a) The expression $(1+2x)^{1/x}$ as $x\to0$.}
For $|x|<1/2$, the base is positive. Its logarithm is $\log(1+2x)/x$, of form $0/0$. The denominator derivative is $1$, and the numerator derivative is $2/(1+2x)\to2$. Thus the logarithm tends to $2$, and the original expression tends to $e^2$.
\step{(b) The expression $(1+2/y)^y$ as $y\to\infty$.}
Set $x=1/y\to0^+$. The expression becomes $(1+2x)^{1/x}$, whose right-hand limit is $e^2$ by part (a). Exercise 30.4 rigorously justifies the reciprocal change of variable.
\step{(c) The expression $(e^x+x)^{1/x}$ as $x\to\infty$.}
The base is positive for $x>0$. Factor out $e^x$ inside its logarithm:
\[
 \frac{\log(e^x+x)}x
 =\frac{x+\log(1+xe^{-x})}x
 =1+\frac{\log(1+xe^{-x})}x.
\]
L'Hospital's rule gives $x/e^x\to0$, so $\log(1+xe^{-x})\to0$. Dividing this term by $x\to\infty$ still gives zero; for example its absolute value is eventually at most $1/x$. The logarithmic limit is $1$, and the original limit is $e$.
\answer{The three limits are $e^2$, $e^2$, and $e$.}

\exercise{30.6}{Recovering limits of a function and its derivative from their sum}
\goal Given $f(x)+f'(x)\to L\in\R$ at positive infinity, prove $f(x)\to L$ and $f'(x)\to0$.
\idea Multiplication by $e^x$ makes the sum appear as a product derivative.
\step{Apply the denominator-infinity version of the rule.}
Write
\[
 f(x)=\frac{f(x)e^x}{e^x}.
\]
Both numerator and denominator are differentiable on $(c,\infty)$, the denominator tends to $+\infty$, and its derivative $e^x$ never vanishes. The derivative quotient is
\[
 \frac{(f(x)e^x)'}{(e^x)'}
 =\frac{e^x(f'(x)+f(x))}{e^x}
 =f'(x)+f(x)\longrightarrow L.
\]
Theorem 30.2 therefore gives $f(x)\to L$. It is unnecessary to know whether the numerator itself has a limit: Ross's denominator-infinity formulation does not require that.
\step{Subtract two now-known finite limits.}
Since $f'(x)=[f(x)+f'(x)]-f(x)$, ordinary limit arithmetic gives $f'(x)\to L-L=0$.
\answer{Both conclusions follow without assuming a limit of $f$ or $f'$ at the outset.}

\exercise{30.7}{Why zeros of the denominator derivative cannot be canceled away}
\goal Analyze $f(x)=x+\cos x\sin x$ and $g(x)=e^{\sin x}f(x)$ and explain the apparent L'Hospital paradox.
\step{(a) Both functions tend to positive infinity.}
Since $|\cos x\sin x|\le1$, we have $f(x)\ge x-1\to\infty$. Also $e^{\sin x}\ge e^{-1}>0$, so for $x>1$,
\[
 g(x)\ge e^{-1}(x-1)\longrightarrow\infty.
\]
\step{(b) Differentiate and factor.}
Using $1=\sin^2x+\cos^2x$,
\[
 f'(x)=1-\sin^2x+\cos^2x=2\cos^2x.
\]
The product and chain rules give
\begin{align*}
 g'(x)&=e^{\sin x}\cos x\,f(x)+e^{\sin x}\,2\cos^2x\\
 &=e^{\sin x}\cos x[2\cos x+f(x)].
\end{align*}
\step{(c) Cancel only at points where cancellation is permitted.}
If $x>3$, then $2\cos x+f(x)\ge-2+x-1=x-3>0$. When also $\cos x\ne0$, the denominator derivative is nonzero and
\[
 \frac{f'(x)}{g'(x)}=
 \frac{2e^{-\sin x}\cos x}{2\cos x+f(x)}.
\]
If $\cos x=0$, however, both derivatives vanish and the left side is undefined. The algebra does not define this original derivative quotient at those points.
\step{(d) Distinguish a continuous extension from the actual quotient.}
The expression on the right in part (c) is defined for every $x>3$, and satisfies
\[
 \left|\frac{2e^{-\sin x}\cos x}{2\cos x+f(x)}\right|
 \le\frac{2e}{x-3}\longrightarrow0.
\]
Nevertheless, $g'(\pi/2+k\pi)=0$ for arbitrarily large integers $k$. Thus there is no tail on which $g'$ is everywhere nonzero, as required by L'Hospital's rule.
For the original functions, $f(x)>0$ on $x>1$, and cancellation of the nonzero $f(x)$ gives
\[
 \frac{f(x)}{g(x)}=e^{-\sin x}.
\]
Along $u_n=\pi/2+2\pi n$ this is $e^{-1}$; along $v_n=3\pi/2+2\pi n$ it is $e$. Both sequences tend to infinity, and the values have different limits. Therefore $f/g$ has no limit.
\answer{The extended derivative expression tends to zero, but the required nonvanishing hypothesis fails. There is no contradiction to L'Hospital's rule.}
